\section{Không có cách so sánh}
\label{sec:ch10-khong-co-cach-so-sanh}

\index{LIME!khả năng dựng lại}%
Khảo sát dành cả một phần thảo luận cho khoảng trống mà mục trước dừng lại ở
đó, và phần ấy mở bằng một câu chẩn đoán: dù LIME phổ biến đến thế, một chuẩn thực
hành cho việc nghiên cứu và đánh giá vẫn chưa được thiết lập, và chính sự thiếu
chuẩn ấy cản trở tính chặt chẽ khoa học của các kỹ thuật dựa trên
LIME~\autocite{p22whichlime}. Ba quan sát cụ thể đứng sau câu đó.

Quan sát thứ nhất là về mã nguồn. Các tác giả ghi nhận một sự thiếu vắng mã
nguồn nghiêm trọng, ở mức 50\%, và nói thẳng rằng nó khiến việc kiểm chứng đóng
góp của các tác giả trở nên khó khăn~\autocite{p22whichlime}. Con số ấy có hai
giới hạn cần nêu. Nó là con số của riêng tập kỹ thuật mà khảo sát gom được,
không phải một cuộc điều tra cả lĩnh vực; và khảo sát dẫn nó về trang web đồng
hành chứ không về một bảng nào trong bài, nên phần thân không nói rõ 50\% ấy là
tỉ lệ thiếu hay tỉ lệ có. Đọc theo câu văn của bài, tức đọc \enquote{thiếu vắng
mã nguồn ở mức 50\%}, thì một nửa số biến thể trong
mục~\ref{sec:ch10-vuon-bien-the} không ai chạy lại được.

Quan sát thứ hai là hệ quả trực tiếp của quan sát thứ nhất, và nó là quan sát
nặng nhất. Vì không có mã nguồn để đối chiếu, khảo sát ghi nhận rằng nhiều bài
mở đầu bằng cách nêu một nhóm vấn đề đã biết của LIME rồi so bản sửa của mình
với đúng một thứ: bản LIME chuẩn. Khảo sát gọi lối làm ấy là bỏ qua các công
trình liên quan và nói nó làm giảm chất lượng nghiên cứu. Những bài có so với
biến thể trước thì thường so với S-LIME hoặc BayLIME, không phải vì hai cái đó
là mốc đúng, mà vì hai cái đó có mã nguồn~\autocite{p22whichlime}.

Hệ quả của lối làm ấy là không có một quan hệ hơn kém nào giữa các biến thể
với nhau. Nếu mỗi biến thể chỉ được so
với gốc, thì sau 48 kỹ thuật ta biết 48 mệnh đề dạng \enquote{biến thể này
tốt hơn LIME chuẩn ở nhóm vấn đề kia}, và không biết một mệnh đề nào so hai
biến thể với nhau. Câu hỏi trong tựa đề khảo sát vì thế vẫn nguyên ở cuối bài,
và khảo sát không đưa ra đáp án nào cho nó.

Quan sát thứ ba là về chỉ số. Khảo sát nêu một vấn đề mà nó gọi là phổ biến,
và gọi là phổ biến trong cả cộng đồng XAI chứ không riêng nhánh LIME: các thước
đo dùng để đánh giá được chọn sao cho chúng xác nhận đóng góp đã tuyên
bố~\autocite{p22whichlime}. Với các đánh giá định tính bằng khảo sát người dùng,
các tác giả ghi nhận không có thủ tục chuẩn nào, cỡ mẫu khác nhau và nhiệm vụ
giao cho người tham gia cũng khác nhau.

Ba quan sát ấy khép thành một vòng, và hình~\ref{fig:ch10-vong-bien-the} vẽ nó.
Nét chấm ở đó có nghĩa là thứ đang thiếu, khác với nét đứt của
hình~\ref{fig:ch08-vong-lap-do}, vốn đánh dấu một lựa chọn để ngỏ.
Một nhóm vấn đề được nêu, một biến thể ra đời để sửa nhóm ấy, rồi câu hỏi biến
thể mới có hơn các biến thể trước hay không cần hai thứ để trả lời: mã nguồn của
các biến thể trước, và một chỉ số mà cả hai bên cùng chấp nhận. Mã nguồn thì
một nửa không có. Chỉ số thì chưa cái nào được kiểm định:
mục~\ref{sec:ch10-mot-tu-khong-xuat-hien} cho thấy khảo sát không có thủ tục
tính nào, và chương~\ref{ch:chi-so-khong-do-do-trung-thuc} cho thấy ngay các
chỉ số có thủ tục cũng chưa được đối chiếu với ground truth. Vì câu hỏi không trả
lời được, việc còn lại mà một tác giả làm được là nêu một nhóm vấn đề khác và
viết biến thể tiếp theo.

\begin{figure}[htbp]
  \centering
  \begin{tikzpicture}[
  every node/.style={font=\footnotesize, align=center},
  box/.style={draw=black!65, rounded corners=2pt, fill=black!5,
    minimum width=30mm, minimum height=11mm},
  askbox/.style={box, draw=black!85, line width=1pt, fill=black!12},
  gapbox/.style={draw=black!45, dotted, thick, rounded corners=2pt, fill=white,
    minimum width=30mm, minimum height=11mm},
  flow/.style={-{Stealth[length=2mm]}, thick, black!70},
  blocked/.style={-{Stealth[length=2mm]}, thick, black!45, dotted},
  note/.style={font=\scriptsize, align=center, black!55, text width=30mm}
]
  \node[box] (issue) at (0,2.4) {Một nhóm vấn đề\\của LIME};
  \node[box] (variant) at (5.8,2.4) {Biến thể mới\\sửa nhóm ấy};
  \node[askbox, minimum width=32mm] (ask) at (5.8,0)
    {Có hơn các biến thể\\trước hay không?};

  \node[gapbox] (metric) at (3.0,-2.4) {Thước đo chung:\\chưa được kiểm chứng};
  \node[gapbox] (code) at (8.6,-2.4) {Mã nguồn:\\một nửa không có};

  \draw[flow] (issue) -- (variant);
  \draw[flow] (variant) -- (ask);
  \draw[blocked] (metric) -- (ask);
  \draw[blocked] (code) -- (ask);
  \draw[blocked] (ask.west) -| (issue.south);

  \node[note, text width=30mm] at (1.9,0.75)
    {không kết luận được,\\nên quay lại viết\\biến thể tiếp theo};
\end{tikzpicture}

  \caption{Ô viền đậm là chỗ duy nhất trong vòng có thể loại bỏ một biến thể, và
  hai đầu vào mà nó cần đều là nét chấm, tức đang thiếu. Bản gộp của sách, dựng
  từ phần thảo luận của khảo sát~\autocite{p22whichlime}.}
  \label{fig:ch10-vong-bien-the}
\end{figure}

Vòng ấy chạy được mà không cần ai gian dối, và đó là điều làm nó khó chữa. Mỗi
bước trong vòng đều hợp lý xét riêng: nêu một hạn chế có thật, đề xuất một cách
sửa, so với bản gốc để cho thấy cách sửa có tác dụng, rồi công bố. Cái sai nằm ở
chỗ vòng không có bước nào loại bỏ được một biến thể; khảo sát không nói câu
sau, nhưng sách đọc từ đó rằng qua chín năm số biến thể chỉ tăng mà chưa cái
nào bị loại.
